\section*{Why the tempting dual of the first head was not used}

The original head is $(\exists x_D\,\forall x_D H)\to H$.
The tempting replacement $H\to\forall x_D\,\exists x_D H$ cannot be
justified by silently assuming ordinary falsity for the distinguished atom 1.
The original head is involved in deriving that explosion interface.

A separating valuation makes codes 0 and 1 true, interprets implication by
the Boolean truth table, interprets each universal operator as the identity
on truth values, interprets $E(a,b)$ as true exactly when $a=b$, and interprets
$M(a,b)$ as true exactly when $a=b=0$. Negation is the defined implication
to code 1; it is therefore always true under this valuation. The existential
operator, defined using this negation and the universal operator, is always
true as well. The dual first head, remaining elementary heads, diagonal
clause and guarded gate are true; modus ponens and theorem generalization
preserve truth. The old first head is false when $H=M(1,1)$, since its
existential antecedent is true and its consequent is false.

Consequently this proposed change is not conservative over the elementary
calculus by the asserted duality argument. It was rejected before any table
was accepted. The example is a warning about proof dependency, not a claim
that every possible consistency calibration of a dual-head system is impossible.
